\section{代数与几何的等价}\label{sec:equivalence-algebra-geometry}

结合本章与上一章的结果, 现在证明从 $X$ 到 $K[X]$ 的转换确实是仿射簇与有限生成既约 $K$-代数之间的一部词典, 其中两边的对象同构时, 就把它们视为\textquotedbl{}相同\textquotedbl{}. \emph{同构类}是同构关系下的等价类, 而同构在 $K$-代数, 仿射簇或更一般的任意范畴中都是等价关系. 用同构类来表述, 已经证明的结果导出如下定理.

\begin{Theorem}{代数与几何的等价}\leavevmode\label{thm:equivalence-algebra-geometry}
  对应 $X\mapsto K[X]$ 诱导双射
  \[
    \{\text{仿射簇的同构类}\}
     \longrightarrow
    \{\text{有限生成既约 }K\text{-代数的同构类}\}.
  \]
\end{Theorem}

\begin{Proof}
由 Proposition~\ref{prop:coordinate-rings-as-finitely-generated-reduced-algebras}, 每个坐标环都是有限生成既约 $K$-代数, 所以可将对应 $X\to K[X]$ 看作函数
\[
  \{\text{仿射簇}\}\longrightarrow
  \{\text{有限生成既约 }K\text{-代数的同构类}\}.
\]
\BookPairedNote
  {为证明这个函数在仿射簇的同构类上良定义, 注意只要 $X\cong Y$, 就有 $K[X]\cong K[Y]$; 这是 Corollary~\ref{cor:coordinate-rings-detect-isomorphisms} 的一个方向. 为证明它在同构类上是单射, 注意只要 $K[X]\cong K[Y]$, 就有 $X\cong Y$; 这是 Corollary~\ref{cor:coordinate-rings-detect-isomorphisms} 的另一个方向.}
  {用范畴论的语言来说, Theorem~\ref{thm:equivalence-algebra-geometry} 的双射反映了仿射簇与有限生成既约 $K$-代数之间的\textquotedbl{}范畴等价\textquotedbl{}.}
\noindent 最后, 为证明满射性, 注意每个有限生成既约 $K$-代数都是某个仿射簇的坐标环 (Proposition~\ref{prop:coordinate-rings-as-finitely-generated-reduced-algebras}).
\end{Proof}

不那么严格地说, Theorem~\ref{thm:equivalence-algebra-geometry} 断言仿射簇 $X$ 的所有几何信息都编码在 $K$-代数 $K[X]$ 中. 不过, 或许应该说得更仔细一些: $X$ 所处的具体仿射空间 $\mathbb{A}^{n}$ 也是关于 $X$ 的\textquotedbl{}几何信息\textquotedbl{}, 但这项信息无法从 $K[X]$ 的同构类恢复. 例如, $\mathcal{V}(y-x^{2})\subseteq\mathbb{A}^{2}$ 与 $\mathcal{V}(y-x^{2},z)\subseteq\mathbb{A}^{3}$ 的坐标环在同一个同构类中, 尽管这两个仿射簇处于不同的环境仿射空间.

之所以无法从 $K[X]$ 的同构类恢复 $X$ 的环境仿射空间, 是因为同构的仿射簇可以处于不同的仿射空间中. 仿射簇, 或者任何数学对象, 在同构下不保持的性质可以看作一种\textquotedbl{}偶然\textquotedbl{}性质, 即依赖于某个外在选择的性质. 相比之下, 同构所保持的性质关乎对象的\textquotedbl{}本质\textquotedbl{}. 为把这些想法说准确, 以下定义很有用.

\begin{Definition}{内在性质与外在性质}\leavevmode\label{def:intrinsic-extrinsic-property}
  设 $C$ 为一组数学对象, 其上有一个称为同构的等价关系. 若每当两个对象同构时, 其中一个具有性质 $P$ 当且仅当另一个具有性质 $P$, 则称 $C$ 中对象的性质 $P$ 是\textbf{内在的}. 不是内在的性质称为\textbf{外在的}.
\end{Definition}

例如, 成为 $\mathbb{A}^{2}$ 中的零点集是所有仿射簇上的一个外在性质, 因为 $\mathcal{V}(y-x^{2})\subseteq\mathbb{A}^{2}$ 有此性质, 而同构的簇 $\mathcal{V}(y-x^{2},z)\subseteq\mathbb{A}^{3}$ 没有此性质. 下面再给出一些代数与几何中的内在性质和外在性质的例子.

\begin{Example}{环的内在性质}\leavevmode\label{ex:intrinsic-ring-properties}
  既约性是环的一个内在性质. 为证明这一点, 设 $\phi:R\to S$ 是环同构; 必须说明 $R$ 既约当且仅当 $S$ 既约. 假设 $R$ 不是既约的. 则存在 $a\in R$ 满足 $a\ne0$, 且对某个 $m\geq1$ 有 $a^{m}=0$. 由环同构的基本性质可知 $\phi(a)\ne0$, 而
  \[
    \phi(a)^{m}=\phi(a^{m})=\phi(0)=0.
  \]
  因此 $\phi(a)$ 是非零幂零元, 说明 $S$ 不是既约的. 将同一证明用于 $\phi^{-1}:S\to R$ 则可说明, 若 $S$ 不是既约的, 那么 $R$ 也不是既约的.

  类似的论证说明, 是整环, UFD, PID 或域都是环的内在性质 (Exercise~\ref{ex:prove-intrinsic-ring-properties}).
\end{Example}

\begin{Example}{生成元个数是外在的}\leavevmode\label{ex:number-generators-extrinsic}
  \strut\par\nopagebreak[4]\noindent
  每个有限生成的 $K$-代数都同构于某个商 $K[x_{1},\ldots,x_{n}]/I$, 但生成元的个数 $n$ 是外在的. 例如, $K$-代数
  \[
    \frac{K[x,y]}{\langle y-x^{2}\rangle}
    \qquad\text{与}\qquad
    \frac{K[x,y,z]}{\langle y-x^{2},z\rangle}
  \]
  是同构的, 尽管第一个用两个生成元, 第二个用三个生成元. 另一方面, 生成元的最少个数是内在性质.
\end{Example}

\begin{Example}{不可约性是内在的}\leavevmode\label{ex:irreducibility-is-intrinsic}
  是不可约仿射簇这一性质是内在的. 这可以直接使用多项式映射和不可约性的定义证明 (Exercise~\ref{ex:prove-irreducibility-intrinsic}), 也可以通过几何与代数之间的词典证明. 为此, 设 $X\subseteq\mathbb{A}^{m}$ 和 $Y\subseteq\mathbb{A}^{n}$ 是仿射簇, 且 $X\cong Y$. 那么
  \[
    \begin{aligned}
      X\text{ 不可约}
        &\Longleftrightarrow \mathcal{I}(X)\text{ 是素理想}
          &&\text{(Proposition~\ref{prop:irreducible-iff-vanishing-ideal-prime})}\\
        &\Longleftrightarrow K[X]\text{ 是整环}
          &&\text{(Propositions~\ref{prop:coordinate-ring-as-quotient-k-algebra} 和 \ref{prop:prime-maximal-quotient-characterization}).}
    \end{aligned}
  \]
  由假设 $X\cong Y$, 于是 $K[X]\cong K[Y]$ (Corollary~\ref{cor:coordinate-rings-detect-isomorphisms}). 由于是整环是 $K$-代数的内在性质, 所以 $X$ 不可约当且仅当 $Y$ 不可约. 因此不可约性是一个内在性质.
\end{Example}

一类数学对象 $C$ 的内在性质 $P$, 可以看成同构类集合的一个子集, 由其中有一个 (从而每一个) 代表具有性质 $P$ 的那些同构类组成. 特别地, 仿射簇的一项内在性质, 比如不可约性, 可以看成仿射簇的同构类集合的一个子集. 通过 Theorem~\ref{thm:equivalence-algebra-geometry} 中的双射, 可以把它等同于有限生成既约 $K$-代数的同构类集合的一个子集, 从而也可以看成有限生成既约 $K$-代数的一项内在性质. 它是哪一项代数性质? 就不可约性而言, 答案是: 仿射簇不可约对应于有限生成既约 $K$-代数是整环.

更一般地, 现在可以问一个很宽泛的问题: 给定一项适用于仿射簇的内在几何性质或构造, 它在 $K$-代数范畴中怎样体现? 反过来, 给定一项适用于 $K$-代数的内在代数性质或构造, 它在仿射簇范畴中怎样体现? 这些问题有助于理解代数, 也有助于理解几何; 这类现象的具体例子构成了后续代数几何的主线.

\BookExerciseHeading
\begin{BookExercises}
  \item\label{ex:prove-intrinsic-ring-properties} 证明下列都是环的内在性质:
    \begin{enumerate}[label=(\alph*)]
      \item 是整环;
      \item 是域;
      \item 是主理想整环;
      \item 是唯一分解整环.
    \end{enumerate}
  \item\label{ex:prove-irreducibility-intrinsic} 设 $X\subseteq\mathbb{A}^{m}$ 和 $Y\subseteq\mathbb{A}^{n}$ 为仿射簇, $F:X\to Y$ 为同构.
    \begin{enumerate}[label=(\alph*)]
      \item 证明若子集 $X_{1}\subseteq X$ 是 $\mathbb{A}^{m}$ 中的仿射簇, 则 $F(X_{1})$ 是 $\mathbb{A}^{n}$ 中的仿射簇.
      \item 直接根据 Definition~\ref{def:reducible-and-irreducible-affine-varieties} 证明不可约性是内在的.
    \end{enumerate}
  \item\label{ex:number-irreducible-components-intrinsic} 证明仿射簇的不可约分支个数是一个内在性质. 描述有限生成既约 $K$-代数中与之对应的性质.
  \item\label{ex:finiteness-of-variety-intrinsic} 证明是有限集是仿射簇的一个内在性质. 描述有限生成既约 $K$-代数中与之对应的性质.
  \item\label{ex:isomorphism-classes-of-morphisms} 若存在同构 $G_{1}:A_{1}\to A_{2}$ 和 $G_{2}:B_{1}\to B_{2}$ 使
    \[
      F_{1}=G_{2}^{-1}\circ F_{2}\circ G_{1},
    \]
    则称两个态射 $F_{1}:A_{1}\to B_{1}$ 和 $F_{2}:A_{2}\to B_{2}$ 是同构的.
    \begin{enumerate}[label=(\alph*)]
      \item 证明同构是态射集合上的一个等价关系.
      \item 证明仿射簇之间多项式映射的同构类与有限生成既约 $K$-代数之间 $K$-代数同态的同构类之间有一个双射.
    \end{enumerate}
\end{BookExercises}
